\documentclass[a4paper,10pt]{article}
\usepackage{graphicx} 
\usepackage[english,russian]{babel}
\usepackage[T2A]{fontenc}
\usepackage[utf8]{inputenc}
\usepackage{amsmath} 
\usepackage{amssymb} 
\usepackage{titlesec}
\usepackage{fancyhdr}
\titleformat{\section}[block]{\centering\large\bfseries}{}{0pt}{}
\titlespacing*{\section}{0pt}{12pt}{12pt}
\pagestyle{fancy}
\fancyhf{}
\fancyhead[L]{\large\bfseries §3}  
\fancyhead[C]{\small\bfseries МЕТОД РАЗДЕЛЕНИЯ ПЕРЕМЕННЫХ} 
\fancyhead[R]{\thepage}                                             
\renewcommand{\headrulewidth}{0pt}  
\begin{document}
\thispagestyle{fancy}
\setcounter{page}{103}
\setcounter{section}{2} 
Полученное выше интегральное представление (3*) решения обыкновенного дифференциального уравнения (1*) имеет, как мы убедились, тот же физический смысл, что и формула (59), дающая интегральное представление решения неоднородного уравнения колебаний.

\textbf{5. Общая первая краевая задача.} Рассмотрим общую первую краевую задачу для уравнения колебаний:
\textit{найти решение уравнения}
\[
u_{tt} = a^2 u_{xx} + f(x, t), \quad 0 < x < l, \, t > 0
\tag{45}
\]
\textit{с дополнительными условиями}
\[
\left.
\begin{aligned}
&u(x, 0) = \varphi(x), \\
&u_t(x, 0) = \psi(x),
\end{aligned}
\right\}
\quad 0 \leq x \leq l;
\tag{46}
\]
\[
\left.
\begin{aligned}
&u(0, t) = \mu_1(t), \\
&u(l, t) = \mu_2(t),
\end{aligned}
\right\}
\quad t \geq 0.
\tag{47}
\]
Введем новую неизвестную функцию \(v(x, t)\), полагая:
\[
u(x, t) = U(x, t) + v(x, t),
\]
так что \(v(x, t)\) представляет отклонение функции \(u(x, t)\) от некоторой известной функции \(U(x, t)\).
Эта функция \(v(x, t)\) будет определяться как решение уравнения
\[
v_{tt} = a^2 v_{xx} +\bar f(x, t), \quad \bar f(x, t) = f(x, t) - \bigl[U_{tt} - a^2 U_{xx}\bigr]
\]
с дополнительными условиями
\[
v(x, 0) = \bar{\varphi}(x), \quad \bar{\varphi}(x) = \varphi(x) - U(x, 0),
\]
\[
v_t(x, 0) = \bar{\psi}(x); \quad \bar{\psi}(x) = \psi(x) - U_t(x, 0);
\]
\[
v(0, t) = \bar{u}_1(t), \quad \bar{u}_1(t) = \mu_1(t) - U(0, t),
\]
\[
v(l, t) = \bar{u}_2(t); \quad \bar{u}_2(t) = \mu_2(t) - U(l, t).
\]
Выберем вспомогательную функцию \(U(x, t)\), таким образом, чтобы
\[
\tilde{\mu}_1(t) = 0 \quad \text{и} \quad \tilde{\mu}_2(t) = 0;
\]
для этого достаточно положить
\[
U(x, t) = \mu_1(t) + \frac{x}{l} \bigl[\mu_2(t) - \mu_1(t)\bigr].
\]
Тем самым общая краевая задача для функции \( u(x, t) \) сведена к краевой задаче для функции \( v(x, t) \) при нулевых граничных условиях. Метод решения этой задачи изложен выше (см. п. 4).
\newpage  

\fancyhead[L]{\thepage}                                      % 104 СЛЕВА
\fancyhead[C]{\bfseries УРАВНЕНИЯ ГИПЕРБОЛИЧЕСКОГО ТИПА}     % НАЗВАНИЕ ПО ЦЕНТРУ
\fancyhead[R]{\bfseries ГЛ.~II}                              % ГЛ. II СПРАВА
\thispagestyle{fancy}
\setcounter{page}{104}
\textbf{6. Краевые задачи со стационарными неоднородностями.} Весьма важным классом задач являются краевые задачи со стационарными неоднородностями, когда граничные условия и правая часть уравнения не зависят от времени
\[
u_{tt} = a^2 u_{xx} + f_0(x), \tag{45'}
\]
\[
\left.
\begin{aligned}
&u(x, 0) = \varphi(x), \\
&u_t(x, 0) = \psi(x),
\end{aligned}
\right\}
\tag{46}
\]
\[
\left.
\begin{aligned}
&u(0, t) = u_1, \quad u_1 = \text{const}, \\
&u(l, t) = u_2, \quad u_2 = \text{const},
\end{aligned}
\right\}
\tag{47'}
\]
В этом случае решение естественно искать в виде суммы
\[
u(x, t) = \bar{u}(x) + v(x, t),
\]
где \(\bar{u}(x)\) — стационарное состояние (статический прогиб) струны, определяемое условиями
\[
\begin{aligned}
a^2 \bar{u}''(x) + f_0(x) &= 0, \\
\bar{u}(0) &= u_1, \\
\bar{u}(l) &= u_2,
\end{aligned}
\]
а \(v(x, t)\) — отклонение от стационарного состояния. Нетрудно видеть, что функция \(\bar{u}(x)\) равна
\[
\bar{u}(x) = u_1 + (u_2 - u_1) \frac{x}{l} + \frac{x}{l} \int_0^l d\xi_1 \int_0^{\xi_1} \frac{f_0(\xi_2)}{a^2} d\xi_2 - \int_0^x d\xi_1 \int_0^{\xi_1} \frac{f_0(\xi_2)}{a^2} d\xi_2.
\]
В частности, если \(f_0 = \text{const}\), то
\[
\bar{u}(x) = u_1 + (u_2 - u_1) \frac{x}{l} + \frac{f_0}{2a^2} (l x - x^2).
\]
Функция \(v(x, t)\), очевидно, удовлетворяет однородному уравнению
\[
v_{tt} = a^2 v_{xx}
\]
с однородными граничными условиями
\[
v(0, t) = 0,
\]
\[
v(l, t) = 0,
\]
и начальными условиями
\[
\begin{aligned}
v(x, 0) &= \bar{\varphi}(x), \quad \bar{\varphi}(x) = \varphi(x) - \bar{u}(x), \\
v_t(x, 0) &= \psi(x).
\end{aligned}
\]

Таким образом, \(v\) является решением простейшей краевой задачи, рассмотренной нами в п. 1 настоящего параграфа.

\end{document}